% Version 3 -- revised manuscript
\documentclass[11pt]{article}
\usepackage[a4paper,margin=2.5cm]{geometry}
\usepackage{amsmath,amssymb,bm,mathtools}
\usepackage{booktabs}
\usepackage{microtype}
\usepackage[hidelinks]{hyperref}

\title{\bfseries Modular Localization, Majorana Bilinears, and Near Saturation of the Tsirelson Bound}
\author{J.~G.~A. Carib\'e$^{1,2}$, M.~S. Guimaraes$^{2}$, I.~Roditi$^{3,4}$, and S.~P. Sorella$^{2}$\\[0.6em]
\small $^{1}$UFRRJ, Universidade Federal Rural do Rio de Janeiro, Departamento de F\'isica,\\
\small Zona Rural, BR-465, Km 07, 23890-000, Serop\'edica, Rio de Janeiro, Brazil\\
\small $^{2}$UERJ -- Universidade do Estado do Rio de Janeiro, Instituto de F\'isica -- Departamento de F\'isica Te\'orica,\\
\small Rua S\~ao Francisco Xavier 524, 20550-013, Maracan\~a, Rio de Janeiro, Brazil\\
\small $^{3}$CBPF -- Centro Brasileiro de Pesquisas F\'isicas, Rua Dr. Xavier Sigaud 150, 22290-180, Rio de Janeiro, Brazil\\
\small $^{4}$Institute for Theoretical Physics, ETH Z\"urich, 8093 Z\"urich, Switzerland\\[0.6em]
\small \texttt{joaogcaribe@ufrrj.br}; \texttt{msguimaraes@uerj.br}; \texttt{roditi@cbpf.br}; \texttt{silvio.sorella@fis.uerj.br}}
\date{}

\begin{document}
\maketitle

\begin{abstract}
We investigate Bell--CHSH correlations generated by dichotomic observables constructed from bilinears of smeared Majorana fields. In contrast with constructions based directly on normalized Majorana fields, the observables considered here are even fermionic composites. We formulate their algebraic normalization and anticommutation requirements, construct modularly localized test functions using Tomita--Takesaki modular operators, and implement the required real-orthogonality conditions through a symplectic Gram--Schmidt procedure. The Bell correlators are evaluated through Wick contractions. A finite-dimensional numerical realization based on six Hermite--Gaussian modes, with Gaussian width $\sigma$, localized around modular frequency $\omega=0$ yields, for $\sigma=10^{-3}$,
\[
\mathcal B_{\rm CHSH}=2.8283708483,
\]
corresponding to approximately $99.998\%$ of the Tsirelson value $2\sqrt2$. The four individual correlations simultaneously approach the ideal CHSH pattern.
\end{abstract}

\section{Introduction}

Bell correlations in relativistic quantum field theory provide a sharp meeting point between locality, entanglement, and the operator-algebraic structure of local observables. The two foundational papers of Summers and Werner develop the Bell framework for local quantum field theory and prove maximal vacuum violation for suitable spacelike separated observables in free Bose and Fermi theories~\cite{SummersWernerI,SummersWernerII}. Their complementary-wedge analysis further shows that maximal violation is generic under the standard algebraic hypotheses of quantum field theory~\cite{SummersWernerGeneric}. These results provide the natural point of departure for studying the Tsirelson bound with a different class of local observables.

The present work considers a different class of observables. Rather than using normalized smeared Majorana fields themselves as the dichotomic observables, we construct even operators from \emph{bilinears} of Majorana fields. These operators belong to the bosonic even sector of the fermionic algebra. Their algebraic construction is consequently different, and their composite character also makes them potentially closer to parity-like and bilinear observables encountered in effective descriptions of Majorana degrees of freedom.

We work with a massive Majorana field in $1+1$-dimensional Minkowski spacetime and with the standard right and left wedge geometry. The modular structure of wedge algebras is controlled by the Bisognano--Wichmann theorem~\cite{BisognanoWichmann1975,BisognanoWichmann1976}, which identifies the modular group with Lorentz boosts. This allows the Tomita--Takesaki operators to be written particularly simply in rapidity space and, after Fourier transformation, in modular-frequency space.

A distinguished role is played by modular frequency $\omega=0$. If $\delta$ denotes the modular operator, its spectral value in the $\omega$ representation is
\begin{equation}
 \lambda(\omega)=e^{-2\pi\omega}.
 \label{eq:lambdaomega_intro}
\end{equation}
Hence $\omega=0$ corresponds precisely to the modular spectral value $\lambda=1$. In the local quantum-field-theoretic setting this spectral structure is tied to the type-$\mathrm{III}_1$ character of local von Neumann algebras. The value $\lambda=1$ is fixed by the modular flow and is therefore a distinguished point of the modular spectrum. In the Summers--Werner analysis this point plays a special role in the mechanism leading to maximal Bell correlations. This observation gives the central motivation for concentrating our construction near $\omega=0$: in the Fourier-dual boost representation, $\omega=0$ is exactly the spectral representative of $\lambda=1$. Numerically, we probe this region with Hermite--Gaussian wave packets of Gaussian width $\sigma$ and study the approach to the Tsirelson value as $\sigma$ decreases.

The paper is organized as follows. Section~\ref{sec:bilinears} introduces the massive $1+1$ dimensional Majorana field in rapidity space, defines the bilinear observables, and establishes their Hermiticity and dichotomicity conditions. Section~\ref{sec:modular} reviews the Bisognano--Wichmann realization of Tomita--Takesaki modular theory for wedges, passes from rapidity to modular-frequency space, constructs Alice's and Bob's localized test functions, and implements the required real-orthogonality constraints by a symplectic Gram--Schmidt procedure. Section~\ref{sec:bell} derives the Bell--CHSH correlators using Wick contractions of smeared Majorana fields. Section~\ref{sec:numerics} presents the Hermite--Gaussian finite-basis implementation and the numerical results. Section~\ref{sec:perspective} discusses possible connections with effective Majorana modes in condensed-matter systems and directions for extending the modular construction.

\section{Majorana Bilinears as Dichotomic Observables}
\label{sec:bilinears}

\subsection{Massive Majorana field and one-particle Hilbert space}

We consider a free massive Majorana field in $1+1$ dimensions in wedge regions, following the modular framework developed in Ref.~\cite{CaribeMajorana2026}. On the positive-energy mass shell, momentum is parametrized by the rapidity $\theta\in\mathbb R$ as
\begin{equation}
 p^0(\theta)=m\cosh\theta,
 \qquad
 p^1(\theta)=m\sinh\theta,
 \label{eq:rapidity}
\end{equation}
where units $c=1$ are used. The one-particle Hilbert space is
\begin{equation}
 \mathcal H_1=L^2(\mathbb R,d\theta),
\end{equation}
and for $f,g\in\mathcal H_1$ the Lorentz-invariant one-particle scalar product in rapidity variables is
\begin{equation}
 \langle f,g\rangle
 =\int_{-\infty}^{\infty}d\theta\,
 \overline{f(\theta)}\,g(\theta).
 \label{eq:innerproducttheta}
\end{equation}
The use of rapidity makes Lorentz boosts act as translations of $\theta$, a fact that will be essential for the modular construction below.

Let $\psi(f)$ denote the Majorana field smeared with $f\in\mathcal H_1$. We use the CAR normalization
\begin{equation}
 \{\psi(f),\psi(g)\}=2\,\Re\langle f,g\rangle,
 \label{eq:CAR}
\end{equation}
so that
\begin{equation}
 \psi(f)^2=\|f\|^2.
 \label{eq:psisquare}
\end{equation}

\subsection{Alice's and Bob's bilinears}

For two one-particle vectors $f,h\in\mathcal H_1$, define
\begin{equation}
 A=-i\,\frac{\psi(h)\psi(f)}{\|h\|\,\|f\|},
 \qquad
 A'=-i\,\frac{\psi(h')\psi(f')}{\|h'\|\,\|f'\|}.
 \label{eq:Alice}
\end{equation}
We require
\begin{equation}
 \Re\langle h,f\rangle=0,
 \qquad
 \Re\langle h',f'\rangle=0.
 \label{eq:Aliceorth}
\end{equation}
By Eq.~\eqref{eq:CAR}, these conditions imply that the two Majorana factors inside each bilinear anticommute. Consequently,
\begin{equation}
 A^\dagger=A,
 \qquad
 (A')^\dagger=A',
\end{equation}
and
\begin{equation}
 A^2=(A')^2=\mathbf 1.
 \label{eq:Adichotomic}
\end{equation}
Thus $A$ and $A'$ are Hermitian dichotomic observables.

For Bob we use the notation
\begin{equation}
 B=-i\,\frac{\psi(p)\psi(g)}{\|p\|\,\|g\|},
 \qquad
 B'=-i\,\frac{\psi(p')\psi(g')}{\|p'\|\,\|g'\|},
 \label{eq:Bob}
\end{equation}
with
\begin{equation}
 \Re\langle g,p\rangle=0,
 \qquad
 \Re\langle g',p'\rangle=0.
 \label{eq:Boborth}
\end{equation}
The same argument gives $B^\dagger=B$, $(B')^\dagger=B'$ and $B^2=(B')^2=\mathbf1$. No additional condition requiring, for example, $\psi(g)$ and $\psi(g')$ to anticommute is imposed.

\section{Tomita--Takesaki Modular Localization in Wedges}
\label{sec:modular}

\subsection{Bisognano--Wichmann theorem in rapidity space}

For the right wedge, the Bisognano--Wichmann theorem identifies the modular flow with Lorentz boosts~\cite{BisognanoWichmann1975,BisognanoWichmann1976}. In rapidity space the one-particle boost generator is
\begin{equation}
 k=-i\frac{d}{d\theta},
 \end{equation}
and, with the conventions used here, the modular operator is
\begin{equation}
 \delta=e^{-2\pi k}.
 \label{eq:deltaK}
\end{equation}
Therefore
\begin{equation}
 (\delta^{1/2}f)(\theta)=f(\theta+i\pi).
 \label{eq:deltahalf_theta}
\end{equation}
For the massive Majorana one-particle space we use the modular conjugation convention
\begin{equation}
 (jf)(\theta)=\overline{f(\theta)},
 \label{eq:jtheta}
\end{equation}
and the Tomita--Takesaki antilinear operator is
\begin{equation}
 s=j\delta^{1/2}.
 \label{eq:stheta_def}
\end{equation}
The real standard subspace associated with the wedge is characterized by the fixed-point condition $sf=f$.

\subsection{Dual modular-frequency representation}

We now pass to the Fourier-dual variable $\omega$, which diagonalizes the boost generator. Our convention is
\begin{equation}
 f(\omega)=\frac{1}{\sqrt{2\pi}}
 \int_{-\infty}^{\infty}d\theta\,
 e^{-i\omega\theta}f(\theta),
 \label{eq:Fourier}
\end{equation}
with
\begin{equation}
 \langle f,g\rangle
 =\int_{-\infty}^{\infty}d\omega\,
 \overline{f(\omega)}g(\omega).
\end{equation}
In this representation,
\begin{align}
 (\delta f)(\omega)&=e^{-2\pi\omega}f(\omega),
 \label{eq:deltaomega}\\
 (\delta^{1/2}f)(\omega)&=e^{-\pi\omega}f(\omega),
 \label{eq:deltahalfomega}\\
 (jf)(\omega)&=\overline{f(-\omega)},
 \label{eq:jomega}\\
 (sf)(\omega)&=e^{\pi\omega}\overline{f(-\omega)}.
 \label{eq:somega}
\end{align}
The adjoint Tomita--Takesaki map acts as
\begin{equation}
 (s^\dagger f)(\omega)=e^{-\pi\omega}\overline{f(-\omega)}.
 \label{eq:sdagomega}
\end{equation}
The modular spectral value is $\lambda(\omega)=e^{-2\pi\omega}$; hence $\omega=0$ is exactly the point $\lambda=1$.

\subsection{Modular localization, standard subspaces, and wedge duality}

We briefly make explicit what is meant here by modular localization; see, for example, the review by Guido~\cite{Guido2011}. Let
\begin{equation}
 \mathcal K(W_R)=\{f\in\mathrm{Dom}(s):sf=f\}
\end{equation}
be the real standard one-particle subspace associated with the right wedge. Standardness means
\begin{equation}
 \mathcal K(W_R)\cap i\mathcal K(W_R)=\{0\},
 \qquad
 \overline{\mathcal K(W_R)+i\mathcal K(W_R)}=\mathcal H_1.
\end{equation}
The Tomita--Takesaki data $(s,j,\delta)$ are the one-particle modular objects associated with this standard subspace.

The projector-like real-linear map $1+s$ maps suitable vectors in $\mathrm{Dom}(s)$ into the fixed-point space $\mathcal K(W_R)$ and is therefore the basic localization map used for Alice. The complementary left wedge is described by the corresponding dual standard subspace. In the fermionic theory ordinary Haag duality is replaced by \emph{twisted duality}; with the phase convention adopted below, Bob's vectors are generated from $1+s^\dagger$ together with the required factors of $i$.

With the present CAR convention, right- and twisted-left-localized one-particle vectors are real-orthogonal,
\begin{equation}
 \Re\langle f,g\rangle=0,
 \qquad
 f\in\mathcal K(W_R),\quad
 g\in\mathcal K(W_L).
 \label{eq:wedge-real-orth}
\end{equation}
Consequently the corresponding odd Majorana fields have the required graded locality. Since the observables $A,A',B,B'$ are even bilinears, the twist disappears at the observable level and complementary-wedge bilinears commute. In particular,
\begin{equation}
 [A,B]=[A,B']=[A',B]=[A',B']=0.
 \label{eq:crosscommute-modular}
\end{equation}
This is the modular-localization origin of the Alice--Bob commutativity required by the Bell construction.

\subsection{Alice's modularly localized test functions}

Let $\varphi(\omega)$ be real. Alice's first pair is generated by
\begin{align}
 f&=(1+s)\varphi,
 \label{eq:fconstruct}\\
 f'&=(1+s)i\varphi.
 \label{eq:fpconstruct}
\end{align}
Using Eq.~\eqref{eq:somega}, these become explicitly
\begin{align}
 f(\omega)&=\varphi(\omega)+e^{\pi\omega}\varphi(-\omega),
 \label{eq:fexplicit}\\
 f'(\omega)&=i\left[\varphi(\omega)-e^{\pi\omega}\varphi(-\omega)\right].
 \label{eq:fpexplicit}
\end{align}

For the second Majorana factor we take real functions $\ell,\eta,\ell',\eta'$ and define
\begin{align}
 h&=(1+s)(\ell+i\eta),
 \label{eq:hconstruct}\\
 h'&=(1+s)(\ell'+i\eta').
 \label{eq:hpconstruct}
\end{align}
Thus
\begin{align}
 h(\omega)
 &=\ell(\omega)+i\eta(\omega)
 +e^{\pi\omega}\left[\ell(-\omega)-i\eta(-\omega)\right],
 \label{eq:hexplicit}\\
 h'(\omega)
 &=\ell'(\omega)+i\eta'(\omega)
 +e^{\pi\omega}\left[\ell'(-\omega)-i\eta'(-\omega)\right].
 \label{eq:hpexplicit}
\end{align}
The required conditions are $\Re\langle f,h\rangle=0$ and $\Re\langle f',h'\rangle=0$.

\subsection{Bob, the commutant, and twisted duality}

Bob's vectors are constructed from the dual Tomita--Takesaki map $s^\dagger$. The factors of $i$ are essential: they implement the fermionic twist required by twisted duality. let $\widehat g(\omega)$ be real and define
\begin{align}
 g&=i(1+s^\dagger)\widehat g,
 \label{eq:gconstruct}\\
 g'&=-i(1+s^\dagger)i\widehat g.
 \label{eq:gpconstruct}
\end{align}
Hence
\begin{align}
 g(\omega)&=i\left[\widehat g(\omega)+e^{-\pi\omega}\widehat g(-\omega)\right],
 \label{eq:gexplicit}\\
 g'(\omega)&=\widehat g(\omega)-e^{-\pi\omega}\widehat g(-\omega).
 \label{eq:gpexplicit}
\end{align}

For the second factor of Bob's bilinears, with real $\alpha,\beta,\alpha',\beta'$, we set
\begin{align}
 p&=i(1+s^\dagger)(\alpha+i\beta),
 \label{eq:pconstruct}\\
 p'&=-i(1+s^\dagger)(\alpha'+i\beta').
 \label{eq:ppconstruct}
\end{align}
Equivalently,
\begin{align}
 p(\omega)
 &=i\Big\{\alpha(\omega)+i\beta(\omega)
 +e^{-\pi\omega}[\alpha(-\omega)-i\beta(-\omega)]\Big\},
 \label{eq:pexplicit}\\
 p'(\omega)
 &=-i\Big\{\alpha'(\omega)+i\beta'(\omega)
 +e^{-\pi\omega}[\alpha'(-\omega)-i\beta'(-\omega)]\Big\}.
 \label{eq:pexplicitprime}
\end{align}
The internal Bob constraints are $\Re\langle g,p\rangle=0$ and $\Re\langle g',p'\rangle=0$.

\subsection{Symplectic Gram--Schmidt procedure}

The orthogonality relevant to the fermionic bilinears is the real part of the one-particle scalar product. It is this real bilinear form that enters the CAR anticommutator in Eq.~\eqref{eq:CAR}; accordingly, the Gram--Schmidt step used here is naturally adapted to the real/symplectic structure of the one-particle space rather than to complex Hilbert-space orthogonality.

For two generic vectors $u,v\in\mathcal H_1$, define
\begin{equation}
 v_{\perp}
 =v-\frac{\Re\langle u,v\rangle}{\|u\|^2}u.
 \label{eq:GSgeneric}
\end{equation}
Then
\begin{equation}
 \Re\langle u,v_{\perp}\rangle=0.
 \label{eq:GSgenericorth}
\end{equation}
We refer to this real projection as the symplectic Gram--Schmidt procedure. It is applied independently to the two Majorana factors entering each bilinear.

For example, for the specific Alice pair, one may define, for a suitable $\xi$,
\begin{equation}
 h_{\perp}
 =\xi-\frac{\Re\langle f,\xi\rangle}{\|f\|^2}f,
 \label{eq:GSfh}
\end{equation}
so that
\begin{equation}
 \Re\langle f,h_{\perp}\rangle=0
 \quad\Longrightarrow\quad
 \{\psi(f),\psi(h_{\perp})\}=0.
 \label{eq:GSfhorth}
\end{equation}
The same operation is performed independently for $(f',h')$, $(g,p)$, and $(g',p')$. In the numerical implementation the projection is imposed before the Bell correlators are evaluated.

\section{Bell--CHSH Correlators and Wick Contractions}
\label{sec:bell}

We now return to the four Majorana bilinears $A,A',B,B'$ constructed above. They satisfy the Bell algebra
\begin{align}
 A^2=(A')^2=B^2=(B')^2&=\mathbf 1,\\
 [A,B]=[A,B']=[A',B]=[A',B']&=0,
 \label{eq:Bellalgebra}
\end{align}
where the first line follows from the internal CAR orthogonality of each bilinear and the second from modular localization in complementary wedges and fermionic twisted duality.

The CHSH operator is
\begin{equation}
 \mathcal C=(A+A')B+(A-A')B'.
 \label{eq:CHSHoperator}
\end{equation}
We denote its vacuum expectation value by
\begin{align}
 \mathcal B_{\rm CHSH}
 :=\langle\Omega,\mathcal C\Omega\rangle
 &=\langle AB\rangle+\langle A'B\rangle
 +\langle AB'\rangle-\langle A'B'\rangle.
 \label{eq:CHSH}
\end{align}
Here $\Omega$ denotes the vacuum vector. Thus $\mathcal C$ always denotes the CHSH operator, whereas $\mathcal B_{\rm CHSH}$ denotes the corresponding Bell parameter in the vacuum state $\Omega$.

Each Alice--Bob correlator is a four-point function of \emph{smeared} Majorana fields. For example,
\begin{equation}
 \langle AB\rangle
 =-\frac{\langle
 \psi(h)\psi(f)\psi(p)\psi(g)
 \rangle}
 {\|h\|\,\|f\|\,\|p\|\,\|g\|}.
 \label{eq:ABfourpoint}
\end{equation}
More generally, for $F_1,F_2,F_3,F_4\in\mathcal H_1$, quasifreeness gives the Wick identity
\begin{align}
 &\left\langle
 \psi(F_1)\psi(F_2)\psi(F_3)\psi(F_4)
 \right\rangle
 \nonumber\\
 &\quad=
 \langle F_1,F_2\rangle\langle F_3,F_4\rangle
 -\langle F_1,F_3\rangle\langle F_2,F_4\rangle
 +\langle F_1,F_4\rangle\langle F_2,F_3\rangle,
 \label{eq:wick}
\end{align}
with the two-point convention inherited from the smeared Majorana field. In particular,
\begin{equation}
 \langle AB\rangle
 =-\frac{1}{\|h\|\|f\|\|p\|\|g\|}
 \left\{
 \langle h,f\rangle\langle p,g\rangle
 -\langle h,p\rangle\langle f,g\rangle
 +\langle h,g\rangle\langle f,p\rangle
 \right\}.
 \label{eq:ABexplicit}
\end{equation}
The other three correlations follow by the corresponding primed substitutions.

\section{Hermite--Gaussian Approximation and Numerical Results}
\label{sec:numerics}

\subsection{Finite trial space}

For Gaussian width $\sigma>0$, the numerical trial space is generated by six Hermite--Gaussian modes, $n=0,\ldots,5$,
\begin{equation}
 \chi_n^{(\sigma)}(\omega)
 =e^{-\omega^2/(2\sigma^2)}
 \frac{H_n(\omega/\sigma)}{\sqrt{2^n n!}}.
 \label{eq:HG}
\end{equation}
The modular seed functions introduced in the preceding section are expanded in this finite basis before the localization maps are applied. For example,
\begin{equation}
 f=(1+s)\varphi,\qquad
 \varphi(\omega)=\sum_{n=0}^{5}c_n\,\chi_n^{(\sigma)}(\omega),
 \label{eq:phiHGexpansion}
\end{equation}
with variational coefficients $c_n$. The analogous expansion is made for each of the remaining real seed functions entering $f',h,h',g,g',p,p'$. Their finite sets of coefficients are then optimized numerically, while the modular localization maps and the symplectic Gram--Schmidt constraints are imposed before the Bell correlators are evaluated.

The small-$\sigma$ regime concentrates the spectral support near $\omega=0$, i.e. near the modular spectral value $\lambda=1$.

\subsection{Convergence toward the Tsirelson value}

The optimized values obtained in the six-mode trial space are
\begin{center}
\begin{tabular}{ccc}
\toprule
$\sigma$ & $\mathcal B_{\rm CHSH}$ & $2\sqrt2-\mathcal B_{\rm CHSH}$\\
\midrule
$1.0\times10^{-2}$ & 2.82474588 & $3.68\times10^{-3}$\\
$5.0\times10^{-3}$ & 2.82718767 & $1.24\times10^{-3}$\\
$2.0\times10^{-3}$ & 2.82820789 & $2.19\times10^{-4}$\\
$1.0\times10^{-3}$ & 2.82837085 & $5.63\times10^{-5}$\\
\bottomrule
\end{tabular}
\end{center}
For reference,
\begin{equation}
 2\sqrt2=2.828427124746\ldots .
\end{equation}
At $\sigma=10^{-3}$,
\begin{equation}
 \mathcal B_{\rm CHSH}=2.8283708483,
\end{equation}
so that
\begin{equation}
 \frac{\mathcal B_{\rm CHSH}}{2\sqrt2}\simeq0.9999801.
\end{equation}
The finite six-mode trial space therefore reaches approximately $99.998\%$ of the Tsirelson value.

The four individual correlators at $\sigma=10^{-3}$ are
\begin{center}
\begin{tabular}{ccc}
\toprule
Correlator & Numerical value & Tsirelson target\\
\midrule
$\langle AB\rangle$   & 0.7071076745  & $+1/\sqrt2$\\
$\langle A'B\rangle$  & 0.7070902621  & $+1/\sqrt2$\\
$\langle AB'\rangle$  & 0.7071025976  & $+1/\sqrt2$\\
$\langle A'B'\rangle$ & -0.7070703141 & $-1/\sqrt2$\\
\bottomrule
\end{tabular}
\end{center}
Thus the four entries approach the ideal CHSH pattern
$(1/\sqrt2,1/\sqrt2,1/\sqrt2,-1/\sqrt2)$ individually, providing a stronger diagnostic than the CHSH sum alone.

The $\sigma=10^{-3}$ result is stable under increasing the Gauss--Hermite quadrature from 120 to 180 and 240 points, yielding
\begin{equation}
 \mathcal B_{\rm CHSH}=2.828370848290538
\end{equation}
at the displayed precision. Within this finite numerical experiment, the observed trend is consistent with
\begin{equation}
 \sigma\to0
 \quad\Longrightarrow\quad
 \omega\to0
 \quad\Longrightarrow\quad
 \lambda\to1
 \quad\Longrightarrow\quad
 \mathcal B_{\rm CHSH}\to2\sqrt2.
 \label{eq:trend}
\end{equation}
This is numerical evidence for near saturation and not, by itself, an analytic proof of the limiting value.

\section{Condensed-Matter Perspective and Outlook}
\label{sec:perspective}

The use of even Majorana bilinears suggests a possible conceptual bridge between the modular field-theoretic construction above and effective Majorana degrees of freedom in condensed-matter systems. Majorana zero modes are predicted in topological superconducting settings and have been extensively studied in atomic spin chains on superconducting substrates and related hybrid structures; for a recent comprehensive review, see Ref.~\cite{RachelWiesendanger2025}.

A possible experimental program would seek effective observables playing the roles of $A,A',B,B'$ while preserving the algebraic ingredients responsible for the large CHSH correlations: dichotomicity, controlled anticommutation within each bilinear, sufficiently independent measurement settings, and an effective analogue of the localization structure used in the modular construction. Because the observables studied here are even Majorana bilinears, parity-sensitive measurements and couplings between effective Majorana modes are natural objects to investigate.

This connection is presently a perspective rather than a microscopic experimental proposal. A concrete realization would require identifying experimentally accessible bilinear/parity measurements, mapping the modularly localized one-particle test functions to controllable effective modes, and analyzing finite temperature, quasiparticle poisoning, decoherence, detector locality, and measurement compatibility.

The present work has deliberately focused on a massive Majorana field in $1+1$ dimensions and wedge localization, where the Bisognano--Wichmann theorem gives direct control of the modular data. A natural extension is to investigate massless/chiral Majorana fields on light rays and localization in bounded regions such as double cones. In that broader setting, modular localization methods of the Brunetti--Guido--Longo type may provide the appropriate framework~\cite{BrunettiGuidoLongo1993}. This extension is left for future work.

\section{Conclusions}

We have constructed Hermitian dichotomic observables from modularly localized Majorana bilinears in the wedge geometry of a massive $1+1$-dimensional Majorana field. The construction makes explicit the role of the CAR real-orthogonality conditions, the Tomita--Takesaki modular maps, twisted duality for the complementary wedge, and the real Gram--Schmidt projection used to enforce the bilinear constraints.

The finite six-mode Hermite--Gaussian implementation approaches the Tsirelson value with high numerical precision as the Gaussian width is reduced and the trial functions become concentrated near modular frequency $\omega=0$, corresponding to the distinguished modular spectral value $\lambda=1$. At $\sigma=10^{-3}$ the CHSH value reaches approximately $99.998\%$ of $2\sqrt2$, while all four individual correlators simultaneously approach the ideal Tsirelson pattern.

An analytic control of the $\sigma\to0$ limit remains an important open problem. Further directions include enlarging the trial space, studying the robustness of the optimizing configurations, extending the modular construction beyond wedges, and investigating whether effective condensed-matter Majorana platforms can realize useful analogues of the bilinear observables considered here.

\section*{Acknowledgments}
The authors acknowledge the use of ChatGPT (OpenAI) as an auxiliary tool for language editing and manuscript preparation. All scientific content, calculations, interpretations, and conclusions were checked and remain the sole responsibility of the authors.

\section*{Data Availability}
The numerical data supporting the findings of this study are contained within the article. Additional numerical details can be made available by the authors upon reasonable request.

\appendix

\section{Numerical Implementation and Precision Analysis}
\label{app:numerics}

This appendix summarizes the numerical implementation used for the finite
Hermite--Gaussian realization of the modular Majorana-bilinear Bell--CHSH
construction and records the corresponding precision checks. The aim is to
separate clearly the variational search, quadrature convergence, arbitrary-
precision reevaluation, and enforcement of the real CAR-orthogonality
constraints.

\subsection{Finite Hermite--Gaussian trial space}

Let $\varepsilon$ denote the modular-frequency variable. For a Gaussian width
$\sigma>0$, the numerical seed functions are expanded in the first six
Hermite--Gaussian modes,
\begin{equation}
\Phi_n^{(\sigma)}(\varepsilon)
=
e^{-\varepsilon^2/(2\sigma^2)}
\frac{H_n(\varepsilon/\sigma)}{\sqrt{2^n n!}},
\qquad
n=0,\ldots,5 .
\label{eq:app-HG}
\end{equation}
Each real seed function entering the modular construction is written as
\begin{equation}
q(\varepsilon)
=
\sum_{n=0}^{5} c_n\,\Phi_n^{(\sigma)}(\varepsilon),
\label{eq:app-seed}
\end{equation}
with real variational coefficients $c_n$. The same six-dimensional ansatz is
used independently for the seed functions from which
$f,f',h,h',g,g',p,p'$ are generated.

The small-$\sigma$ regime concentrates the spectral support near
$\varepsilon=0$. In the modular spectral parametrization employed in the main
text, this is the region corresponding to the distinguished modular spectral
value $\lambda=1$.

\subsection{Modular localization before optimization}

For every trial point, the modular localization maps are applied before the Bell
correlators are evaluated. Alice's vectors are generated with $1+s$, whereas
Bob's vectors are generated with $1+s^\dagger$, including the fermionic twist
phases used in the main construction. Thus the variational coefficients belong
to the seed functions, while the objective function is evaluated on the
corresponding modularly localized vectors.

The numerical ordering is therefore
\begin{equation}
\text{seed coefficients}
\;\longrightarrow\;
\text{modularly localized vectors}
\;\longrightarrow\;
\text{CAR projection}
\;\longrightarrow\;
\mathcal B_{\rm CHSH}.
\label{eq:app-ordering}
\end{equation}
In particular, the routine does not optimize unconstrained Bell correlators and
only afterwards test localization.

\subsection{Symplectic Gram--Schmidt projection and CAR constraints}

For a pair of one-particle vectors $u,v$, the relevant fermionic constraint is
\begin{equation}
\Re\langle u,v\rangle=0 .
\label{eq:app-car}
\end{equation}
The numerical projection is the real, or symplectic, Gram--Schmidt step
\begin{equation}
v_\perp
=
v-
\frac{\Re\langle u,v\rangle}{\|u\|^2}\,u .
\label{eq:app-gs}
\end{equation}
It is applied independently to
\begin{equation}
(f,h),\qquad
(f',h'),\qquad
(g,p),\qquad
(g',p').
\label{eq:app-pairs}
\end{equation}

At the final high-precision reevaluation, the residual real-orthogonality
constraints were found to be
\begin{align}
\Re\langle f,h\rangle
&\sim 7.1\times10^{-62},
\\
\Re\langle f',h'\rangle
&\sim 2.8\times10^{-61},
\\
\Re\langle g,p\rangle
&\sim -3.0\times10^{-61},
\\
\Re\langle g',p'\rangle
&\sim 1.1\times10^{-62}.
\label{eq:app-residuals}
\end{align}
Thus the internal CAR constraints remain satisfied essentially at the full
working precision of the final reevaluation.

\subsection{Evaluation of the Bell correlators}

Each Bell correlator is a four-point vacuum expectation value of smeared
Majorana fields. Quasifreeness reduces these quantities to one-particle scalar
products through Wick's theorem,
\begin{align}
\langle
\psi(F_1)\psi(F_2)\psi(F_3)\psi(F_4)
\rangle
&=
\langle F_1,F_2\rangle\langle F_3,F_4\rangle
-
\langle F_1,F_3\rangle\langle F_2,F_4\rangle
\nonumber\\
&\quad+
\langle F_1,F_4\rangle\langle F_2,F_3\rangle .
\label{eq:app-wick}
\end{align}
The normalization factors associated with the Majorana bilinears are then
included explicitly. The Bell parameter is
\begin{equation}
\mathcal B_{\rm CHSH}
=
\langle AB\rangle
+
\langle A'B\rangle
+
\langle AB'\rangle
-
\langle A'B'\rangle .
\label{eq:app-bchsh}
\end{equation}
Once the localized vectors have been constructed, the numerical problem thus
reduces to repeated evaluation of one-dimensional scalar products in
$\varepsilon$.

\subsection{Gauss--Hermite quadrature}

Because the basis functions in Eq.~\eqref{eq:app-HG} carry a Gaussian factor,
Gauss--Hermite quadrature provides a natural discretization of the required
one-dimensional integrals. For $\sigma=10^{-3}$, the final optimized point was
reevaluated without changing the variational coefficients while increasing the
quadrature order. The result is
\begin{table}[h]
\centering
\begin{tabular}{cc}
\hline\hline
Gauss--Hermite order $N$ & $\mathcal B_{\rm CHSH}$ \\
\hline
80  & 2.828406911727749 \\
120 & 2.828406911727749 \\
180 & 2.828406911727749 \\
240 & 2.828406911727749 \\
\hline\hline
\end{tabular}
\caption{Quadrature convergence at $\sigma=10^{-3}$ for the reconstructed
six-mode optimum.}
\label{tab:app-quadrature}
\end{table}
The displayed value is therefore stable to all 15 digits shown already from
$N=80$ through $N=240$. This test controls the discretization of the integral
and should be distinguished from the floating-point precision used in the
variational search.

\subsection{Optimization precision and high-precision certification}

The numerical workflow was carried out in two stages.

First, the variational optimization was performed in ordinary double precision.
This stage is computationally efficient and is used to locate a high-quality
point in the finite-dimensional coefficient space.

Second, the optimized coefficient vector was held fixed and all relevant scalar
products, Gram--Schmidt projections, normalizations, Wick contractions, and the
final CHSH combination were reevaluated using 60-digit arithmetic. The resulting
value is
\begin{equation}
\mathcal B_{\rm CHSH}
=
2.8284069117277489952982694957984961868084126360042 .
\label{eq:app-highprecision}
\end{equation}
The residual imaginary part of the final Bell value is of order $10^{-61}$ and
is therefore numerically zero at the working precision.

It is important to distinguish these two stages. The nonlinear optimization
itself was not performed at 60-digit precision; rather, the optimum was located
in double precision and subsequently reevaluated and numerically certified at
60 digits.

\subsection{Distance from the Tsirelson bound}

The exact Tsirelson value is
\begin{equation}
2\sqrt{2}
=
2.8284271247461900976\ldots .
\label{eq:app-tsirelson}
\end{equation}
For the reconstructed six-mode optimum,
\begin{equation}
2\sqrt{2}-\mathcal B_{\rm CHSH}
\simeq
2.02130184411023\times10^{-5},
\label{eq:app-distance}
\end{equation}
or equivalently
\begin{equation}
\frac{\mathcal B_{\rm CHSH}}{2\sqrt{2}}
\simeq
0.999992853\ldots .
\label{eq:app-ratio}
\end{equation}
Thus the finite six-mode realization reaches approximately $99.9993\%$ of the
Tsirelson bound.

The four individual correlators at this point are
\begin{align}
\langle AB\rangle
&=
0.7069610194038407592,
\\
\langle A'B\rangle
&=
0.7072492746151345647,
\\
\langle AB'\rangle
&=
0.7070410913686358234,
\\
\langle A'B'\rangle
&=
-0.7071555263401378480.
\label{eq:app-correlators}
\end{align}
Each entry is close to the ideal Tsirelson pattern
\begin{equation}
\left(
\frac{1}{\sqrt{2}},
\frac{1}{\sqrt{2}},
\frac{1}{\sqrt{2}},
-\frac{1}{\sqrt{2}}
\right),
\label{eq:app-pattern}
\end{equation}
which provides a stronger diagnostic than the CHSH sum alone.

\subsection{Reproducibility remark}

The variational coefficients associated with the earlier exploratory run quoted
in the main manuscript were not archived. The numerical calculation summarized
here was therefore reconstructed independently from the equations defining the
same six-mode variational problem. The reconstructed optimum is slightly better
than the earlier value.

Accordingly, the result in Eq.~\eqref{eq:app-highprecision} should be understood
as a new optimum of the same finite Hermite--Gaussian trial problem, rather than
as a literal rerun of an archived coefficient vector.

\subsection{Scope of the numerical checks}

The convergence and precision tests reported above establish that, for the
reconstructed finite-dimensional optimum, the Bell value is insensitive to
quadrature refinement and stable under arbitrary-precision reevaluation. They
also verify that the internal CAR constraints are satisfied to extremely high
numerical accuracy.

These checks do not constitute an analytic proof that
\begin{equation}
\mathcal B_{\rm CHSH}\longrightarrow 2\sqrt{2}
\qquad
\text{as}
\qquad
\sigma\longrightarrow0,
\end{equation}
nor do they prove that the six-mode point is the global optimum of the full
infinite-dimensional variational problem. Their appropriate interpretation is
therefore as high-accuracy numerical evidence for near saturation of the
Tsirelson bound within the finite Hermite--Gaussian realization.









\begin{thebibliography}{99}

\bibitem{SummersWernerI}
S.~J. Summers and R.~Werner,
``Bell's inequalities and quantum field theory. I. General setting,''
J. Math. Phys. \textbf{28}, 2440--2447 (1987),
doi:10.1063/1.527733.

\bibitem{SummersWernerII}
S.~J. Summers and R.~Werner,
``Bell's inequalities and quantum field theory. II. Bell's inequalities are maximally violated in the vacuum,''
J. Math. Phys. \textbf{28}, 2448--2456 (1987).

\bibitem{SummersWernerGeneric}
S.~J. Summers and R.~Werner,
``Maximal violation of Bell's inequalities is generic in quantum field theory,''
Commun. Math. Phys. \textbf{110}, 247--259 (1987),
doi:10.1007/BF01207366.

\bibitem{BisognanoWichmann1975}
J.~J. Bisognano and E.~H. Wichmann,
``On the duality condition for a Hermitian scalar field,''
J. Math. Phys. \textbf{16}, 985--1007 (1975).

\bibitem{BisognanoWichmann1976}
J.~J. Bisognano and E.~H. Wichmann,
``On the duality condition for quantum fields,''
J. Math. Phys. \textbf{17}, 303--321 (1976).

\bibitem{Tsirelson1980}
B.~S. Cirel'son,
``Quantum generalizations of Bell's inequality,''
Lett. Math. Phys. \textbf{4}, 93--100 (1980).

\bibitem{Guido2011}
D.~Guido,
``Modular theory for the von Neumann algebras of local quantum physics,''
in \emph{Aspects of Operator Algebras and Applications},
Contemp. Math. \textbf{534}, 97--120 (2011), arXiv:0812.1511.

\bibitem{CaribeMajorana2026}
J.~G.~A. Carib\'e, M.~S. Guimaraes, I.~Roditi and S.~P. Sorella,
``Modular wedge localization, Majorana fields and the Tsirelson limit of the Bell--CHSH inequality,''
arXiv:2605.06224 [hep-th] (2026).

\bibitem{BrunettiGuidoLongo1993}
R.~Brunetti, D.~Guido, and R.~Longo,
``Modular structure and duality in conformal quantum field theory,''
Commun. Math. Phys. \textbf{156}, 201--219 (1993).

\bibitem{RachelWiesendanger2025}
S.~Rachel and R.~Wiesendanger,
``Majorana quasiparticles in atomic spin chains on superconductors,''
Phys. Rep. \textbf{1099}, 1--28 (2025),
doi:10.1016/j.physrep.2024.10.005.

\end{thebibliography}

\end{document}
